\chapter{Introducción}

\section{Planteamiento del problema}

La simulación de sistemas complejos ha surgido como una herramienta indispensable para la comprensión y el análisis de fenómenos en diversas áreas del conocimiento, desde la física y la biología hasta las ciencias sociales y económicas. Tradicionalmente, los autómatas celulares han servido como un modelo matemático fundamental para estudiar estos sistemas mediante la representación de procesos dinámicos. 

Los autómatas celulares son un modelo matemático definido como una 4-tupla que contiene: 

\begin{itemize}
	
	\item Un espacio de celdas en un dominio discreto.
	
	\item Un conjunto de estados posibles.
	
	\item Una vecindad que define las dependencias entre células.
	
	\item Una función de transición que determina cómo evolucionan simultáneamente las células \cite{weisstein2002_ca, ZenilMartinez2024}. 
	
\end{itemize}

Ejemplos clásicos de autómatas celulares incluyen el juego de la vida de John Conway y los 256 autómatas elementales descritos por Stephen Wolfram. Estas herramientas han sido aplicadas en múltiples contextos como la simulación de incendios, modelado de distribuciones poblacionales, simulación de partículas y el análisis de ciclos climáticos \cite{darg}. No obstante, una limitación crítica de los modelos convencionales reside en la rigidez de sus funciones de transición; estas deben definirse de manera estática y anticipada, lo que obliga a conocer de antemano la regla del sistema que se quiere simular. Esta limitación motiva la búsqueda de nuevas metodologías capaces de incorporar mayor flexibilidad y adaptabilidad en la simulación de sistemas dinámicos. 

En este contexto, se busca la integración de redes neuronales en el marco de los autómatas celulares. Las redes neuronales, inspiradas en el funcionamiento del cerebro y su capacidad de aprendizaje, ofrecen la posibilidad de aprender las reglas de transición a partir de datos. Esto permite que la regla del autómata celular se obtenga mediante entrenamiento en lugar de definirse a mano \cite{RN}. 

El presente proyecto de investigación propone superar las limitaciones de los autómatas celulares mediante la integración de redes neuronales en la definición de las reglas de transición, dando origen a los \textit{autómatas celulares neuronales}. En este enfoque, las reglas de actualización dejan de definirse a mano y se aprenden a partir de los datos, lo que permite obtener reglas para comportamientos complejos cuya regla se desconoce \cite{matich2001_ann, mordvintsev2020_nca}. Esta aproximación no solo busca ampliar el campo de aplicación de los autómatas celulares, sino que también busca abrir nuevas oportunidades para modelar y simular fenómenos de formas realistas.  

Dadas las infinitas configuraciones de autómatas celulares, este proyecto se limita a un espacio bidimensional finito (con condiciones de frontera) y un conjunto finito y numerable de estados. Adicionalmente, se busca experimentar con diferentes vecindades, incluyendo Moore y von Neumann. Aunque el objetivo no es determinar las condiciones para la evolución de las células, sí se busca realizar pruebas bajo diferentes configuraciones que permitan mostrar los resultados sin comprometer su calidad.

Por otro lado, la red neuronal recibe como parámetros los estados de las células vecinas y devuelve el nuevo estado de la célula a la cual se aplique la función de transición. Este proceso es el mismo para todas las células.

Además, la configuración del autómata celular neuronal que se propone difiere de las que se mencionan en el estado del arte. Aquí se experimenta con autómatas celulares neuronales en un espacio discreto bidimensional, cuyas células toman estados finitos, discretos y escalares, y cuya evolución es determinista. Con ello se busca demostrar la capacidad de replicar comportamientos complejos sin necesidad de estados continuos, estados ocultos ni actualizaciones estocásticas (asíncronas).

Por esto surge la pregunta central de esta investigación: ¿cómo pueden los autómatas celulares neuronales simular sistemas complejos en un espacio bidimensional? A partir de este cuestionamiento, surgen las siguientes preguntas de investigación.

\subsection{Preguntas de investigación}

\subsubsection{Preguntas teóricas}

\begin{itemize}
	
	\item ¿Qué es un sistema complejo y qué significa simularlo?
	
	\item ¿Qué es un autómata celular?
	
	\item ¿Qué es una red neuronal?
	
	\item ¿Qué es un autómata celular neuronal?
	
\end{itemize}

\subsubsection{Preguntas experimentales}

\begin{itemize}
	
	\item ¿Puede un autómata celular neuronal inferir de forma exacta la regla de un autómata celular conocido a partir de su evolución?
	
	\item ¿Qué configuración de la red neuronal (capas, neuronas, función de activación, codificación de la entrada y vecindad) se necesita para aprender una regla?
	
	\item ¿Cómo influyen los parámetros de entrenamiento, como la pérdida objetivo y la tasa de aprendizaje, en la regla aprendida?
	
	\item ¿Qué condiciones debe cumplir un sistema real para que un autómata celular neuronal pueda replicarlo a partir de un conjunto de datos?
	
\end{itemize}

\section{Objetivo}

Desarrollar un sistema de simulación de autómatas celulares neuronales bidimensionales que permita inferir, a partir de la evolución de un sistema complejo, la regla local que la rige.

\subsection{Objetivos específicos}

\begin{itemize}
	
	\item Diseñar e implementar un simulador de espacios discretos bidimensionales que sirva de plataforma para estudiar los autómatas celulares neuronales.
	
	\item Entrenar y evaluar el desempeño de los autómatas celulares neuronales bajo diferentes configuraciones.
	
	\item Replicar configuraciones clásicas con comportamientos y soluciones conocidos.
	
	\item Evaluar la capacidad de los autómatas celulares neuronales para replicar sistemas reales a partir de conjuntos de datos.
	
	\item Analizar cómo influyen la configuración de la red y del entrenamiento en la regla aprendida.
	
\end{itemize}

\section{Hipótesis}

Si la función de transición de un autómata celular se define mediante una red neuronal entrenada con la evolución de un sistema, entonces es posible inferir una regla local que reproduzca dicha evolución.

\section{Justificación}

Los autómatas celulares son modelos lo bastante simples para permitir un análisis teórico riguroso, pero lo suficientemente complejos para exhibir fenómenos emergentes. Definir las reglas locales que, aplicadas de forma síncrona, describan comportamientos erráticos o no lineales resulta una tarea compleja \cite{computationalBiology2016}. 

Los modelos tradicionales de autómatas celulares se basan en reglas de transición fijas y ya establecidas. Aunque este enfoque ha sido útil para entender algunos fenómenos, no siempre refleja la complejidad de los sistemas dinámicos que encontramos en la naturaleza, en la sociedad y en la tecnología. 

Cada autómata celular se diseña para un comportamiento particular. Sin embargo, su definición matemática implica la existencia de una infinidad de posibles configuraciones y comportamientos. Por ello, llevar a cabo ajustes manuales hasta obtener el resultado deseado se convierte en un proceso computacionalmente inviable \cite{elYacoubi2021control}. 

Por ello, se usan redes neuronales para definir las reglas de transición, con el fin de crear un modelo de autómata celular neuronal cuya regla se aprenda a partir de la evolución observada de un sistema. Esto no solo amplía el alcance y la aplicabilidad de los autómatas celulares, sino que también abre un novedoso campo de estudio en el que se combinan herramientas de inteligencia artificial y simulación de sistemas complejos. 

La principal innovación de este enfoque reside en reemplazar las funciones de transición estáticas por modelos ajustables. Usando técnicas de aprendizaje automático, se puede entrenar la función de transición con datos representativos. Esto permite obtener la regla de un sistema sin definirla a mano \cite{gilpin2019_nca}. 

La elección de redes neuronales para modelar la función de transición no es arbitraria. Varios estudios recientes han mostrado que las redes neuronales son muy efectivas para aproximar funciones no lineales y entender relaciones complejas en el espacio y el tiempo \cite{mordvintsev2020_nca}. En este trabajo se usa un perceptrón multicapa que recibe los estados de la vecindad de una célula y devuelve su nuevo estado, ya que la función de transición de un autómata celular es local. Estas características las hacen particularmente adecuadas para emular las dinámicas locales y temporales propias de los autómatas celulares. Este proyecto se basa en una línea de investigación establecida que intenta unir la capacidad expresiva de las redes neuronales con la estructura organizada y distribuida de los autómatas celulares. Así, se busca crear un modelo más flexible y aplicable a diferentes situaciones.

Adicionalmente, este proyecto profundiza en técnicas de simulación, modelado matemático y redes neuronales, integrando conocimientos de ciencias de la computación, vida artificial, matemáticas e inteligencia artificial.

La experiencia adquirida durante el desarrollo e implementación de estos modelos abre puertas a futuras investigaciones y proyectos en el campo de la inteligencia artificial. Además, el simulador es de código abierto, lo que contribuye al avance del estado del arte y facilita que otros investigadores puedan extender y aplicar estos conceptos.

En resumen, este proyecto contribuye al desarrollo de la simulación de sistemas complejos mediante un modelo de autómatas celulares neuronales flexible y extensible, que une la teoría de sistemas dinámicos con las técnicas actuales de aprendizaje automático.

\section{Alcances y limitaciones}

\textbf{Alcances}

\begin{itemize}
	
	\item El estudio se limita a la simulación en espacios bidimensionales finitos con condiciones de frontera establecidas en una red regular definida.
	
	\item Se emplea aprendizaje supervisado para definir la función de transición mediante un perceptrón multicapa.
	
	\item La regla aprendida es determinista y se aplica de forma síncrona a todas las células.
	
	\item El simulador admite hasta 20 estados y vecindades semitotalísticas de hasta $7 \times 7$ células.
	
	\item Se analizan comportamientos de patrones conocidos y emergentes bajo distintas configuraciones de estados y vecindades.
	
	\item Se replican sistemas reales a partir de conjuntos de datos de su evolución.
	
\end{itemize}

\textbf{Limitaciones}

\begin{itemize}
	
	\item El simulador y los experimentos no contemplan la generalización del modelo a diferentes dimensionalidades o topologías de red.
	
	\item El análisis se centra en la simulación de sistemas complejos, sin abordar la optimización computacional del entrenamiento.
	
	\item Los resultados se interpretan únicamente dentro del contexto de los parámetros y condiciones experimentales definidos.
	
\end{itemize}

